
\section{同伦幺元的构造}
\label{sec:7-3}
\liftIndex{I280}%

\subsection{结果陈述与证明思路}
\label{subsec:7-3-1}
本节旨在证明如下结果。

\begin{theorem}{7.3.1}\label{thm:7-3-1}
设 $\mathfrak m$ 为第~\ref{sec:7-2} 节
利用所构造的基本链定义的算子，
$\mathcal X_0(L)$ 为该节构造的
奇异单形可数集。
则可找到包含 $\mathcal X_0(L)$ 的
奇异单形可数集 $\mathcal X(L)$，
并将 $\mathfrak m$ 延拓到其上，
使之定义一个滤过 $A_\infty$ 代数，
而 $\mathbf e=PD([L])$，即 $L$ 的
基本闭链的庞加莱对偶，是同伦幺元。
\end{theorem}

（注意，定理~\ref{thm:7-3-1} 中的
$\mathcal X_0(L)$
与第~\ref{subsec:7-2-4} 小节的
$\mathcal X_0(L)$ 不同。
换言之，这里不再区分第~\ref{subsec:7-2-4} 小节
所用的“代”。）

开始证明之前，先说明证明的主要思路，
并指出直接实施时的一处困难。
为此，考虑以下类型的遗忘映射
\liftIndex{I281}。
令 $\mathcal M_{k+1}(\beta)$ 为同伦类
$\beta$ 中、边界上有 $k+1$ 个标记点的
伪全纯圆盘模空间，
并将其视为带有仓西结构的空间。
给定
$\vec a=(a_1,\ldots,a_{|\vec a|})
\in\{1,\ldots,k\}^{|\vec a|}$，
其中 $a_1<\cdots<a_{|\vec a|}$，
令
\[
\operatorname{forget}_{\vec a}\colon
\mathcal M^{\mathrm{main}}_{k+1}(\beta)
\longrightarrow
\mathcal M^{\mathrm{main}}_{k+1-|\vec a|}(\beta)
\]
表示忘掉编号为
$a_1,\ldots,a_{|\vec a|}$ 的标记点的映射。
按构造，仓西结构将与这些遗忘映射相容。
（相容性的确切含义见引理 7.3.8。）

若 $L$ 的光滑奇异单形
$P=(|P|,f)$ 可写成
$f=f'\circ\pi$，
其中 $P'=(|P'|,f')$ 也是光滑奇异单形，
$\dim|P'|=\dim|P|-1$，
而 $\pi\colon|P|\to|P'|$ 是单纯映射，
便称 $P$ 为\emph{退化单形}\liftIndex{I332}。
退化单形的线性组合称为\emph{退化链}。

\begin{lemma}{7.3.2}\label{lem:7-3-2}
若 $\mathcal M^{\mathrm{main}}_{k+1}(\beta)$
上有一族与遗忘映射
$\operatorname{forget}_{\vec a}$ 相容的
多值截面 $\mathfrak s_{k+1,\beta}$，
并有满足全部所需横截性条件的
奇异单形可数集 $\mathcal X(L)$，
则
\begin{equation}
\mathfrak m_{k+1,\beta}
(P_1,\ldots,P_i,L,P_{i+1},\ldots,P_k)
\begin{cases}
\equiv 0,
  &(k,\beta)\ne(1,\beta_0),\\
\equiv\pm P_1,
  &(k,\beta)=(1,\beta_0),
\end{cases}
\tag{7.3.3}\label{eq:7-3-3}
\end{equation}
其中 $\equiv$ 表示模退化链相等。
\end{lemma}

\begin{samepage}
\begin{remark}{7.3.4}\label{rem:7-3-4}
\begin{enumerate}[label=\textup{(\arabic*)},leftmargin=2.8em]
\item 引理~\ref{lem:7-3-2} 表明，
在滤过 $A_\infty$ 代数中，
$PD([L])$ 模退化链子复形是幺元。
\item 实际上，这是一个假设性的引理，
因为难以构造满足引理~\ref{lem:7-3-2}
之假设的结构。
\end{enumerate}
\end{remark}
\end{samepage}

\begin{proof}[引理~\ref{lem:7-3-2} 的证明]
置
\[
\vec P^+=(P_1,\ldots,P_i,L,P_{i+1},\ldots,P_k),
\qquad
\vec P=(P_1,\ldots,P_i,P_{i+1},\ldots,P_k),
\]
并使用模空间
\[
\begin{aligned}
\mathcal M^{\mathrm{main}}_{k+2}(\beta;\vec P^+)
&=\mathcal M^{\mathrm{main}}_{k+2}(\beta)
\mathbin{\times_{L^{k+1}}}
(P_1\times\cdots\times P_i\times L
\times P_{i+1}\times\cdots\times P_k),\\
\mathcal M^{\mathrm{main}}_{k+1}(\beta;\vec P)
&=\mathcal M^{\mathrm{main}}_{k+1}(\beta)
\mathbin{\times_{L^k}}
(P_1\times\cdots\times P_i
\times P_{i+1}\times\cdots\times P_k).
\end{aligned}
\]
作为集合，有
\[
\operatorname{ev}_0
\bigl(\mathcal M^{\mathrm{main}}_{k+2}
(\beta;\vec P^+)\bigr)
=\operatorname{ev}_0
\bigl(\mathcal M^{\mathrm{main}}_{k+1}
(\beta;\vec P)\bigr).
\]
另一方面，当 $(k,\beta)\ne(1,\beta_0)$ 时，
有
\begin{equation}
\dim\mathcal M^{\mathrm{main}}_{k+2}(\beta;\vec P^+)
=\dim\mathcal M^{\mathrm{main}}_{k+1}(\beta;\vec P)+1.
\tag{7.3.5}\label{eq:7-3-5}
\end{equation}
这里的 $\dim$ 是\emph{虚维数}，
即仓西结构意义下的维数。
$\mathcal M^{\mathrm{main}}_{k+2}(\beta;\vec P^+)$
多出的一维，对应于沿 $\partial\Sigma$
参数化第 $i+1$ 个标记点位置的参数。
在引理~\ref{lem:7-3-2} 的假设下，
可以找到
$\mathcal M^{\mathrm{main}}_{k+2}(\beta;\vec P^+)$
的一个扰动（多值截面）
$\mathfrak s_{k+2,\beta,\vec P^+}$，
使之与
$\mathcal M^{\mathrm{main}}_{k+1}(\beta;\vec P)$
上的扰动相容。
（$\mathfrak s_{k+2,\beta,\vec P^+}$
中省去了表示装饰的后缀，
因为它已由 $\vec P^+$ 自动决定。）
由 (7.3.5)，可以选取
$\mathcal M^{\mathrm{main}}_{k+2}
(\beta;\vec P^+)^{\mathfrak s_{k+2,\beta,\vec P^+}}$
的三角剖分，使其虚基本链成为退化链。
引理~\ref{lem:7-3-2} 得证。
\end{proof}

如前所述，由于上一节一直讨论的
横截性问题，似乎难以找到满足
引理~\ref{lem:7-3-2} 假设的
模空间 $\mathcal M^{\mathrm{main}}_{k+1}(\beta)$
的扰动。这就是我们构造同伦幺元
而非严格幺元的原因。
还需说明，定理~\ref{thm:7-3-1} 的证明
并不使用退化链的概念，
因为不容易构造积结构 $\mathfrak m_k$，
使退化链之间的积仍是退化链。
本节末尾将通过更仔细的论证，
把条件 7.3.14 中“模退化链相等”
替换为作为奇异链的严格相等。
为克服这一困难，我们考虑模空间上的
两个（不同的）多值截面：
一个与遗忘映射相容，
另一个是第~\ref{sec:7-2} 节所用的截面。
随后取这两个多值截面之间的同伦，
以得到定义~\ref{def:3-3-2}
所要求的算子 $\mathfrak h_k$。
事实上，需要的是一族以 $[0,1]^k$
为参数的多值截面，
而不只是一个同伦。
